\chapter*{Introduction}
\addcontentsline{toc}{chapter}{Introduction}

At two in the morning the run finally works. The file lands on the instrument computer, and before it can go anywhere you have to say what it is. Which folder. What to call it. Which of today's ten files this one turned out to be. You are being asked to write down what a thing is before anyone knows what it means, and you are being asked at the hour you are least able to answer. So you do what everyone does. You call it \texttt{final\_v3\_actually.dat} and you go home.

A year later nobody can say which \labSampleWord{} it came from, or who took it, or whether the copy on this drive is the same bytes as the copy on that one. The data is not lost. It is worse than lost. It is sitting right there, and you cannot trust it.

The usual remedy is to design a schema first: decide what the entities are, what fields they have, how they connect, and then collect. That works when you already know what you are collecting. It fails in the case that matters, which is when the useful connections are the ones nobody anticipated. And it fails hardest at the moment of saving, where the cost of the decision is highest and the information available to make it is lowest. It also asks the wrong person to do the job. It asks whoever is at the instrument at two in the morning to invent a categorization before they are allowed to save anything. They will not do it, and they are right not to. The task is impossible as posed.

This system runs the operations in the other order.

Save the bytes the instant they exist, under a name that takes no decision at all, because it is computed from the bytes themselves. A standard hash function over a file returns a short string that is the same on every computer and in every decade. It is not a catalog number that somebody has to mint and everybody has to agree on. It is a fingerprint. Two people who capture the same file arrive at the same name without talking to each other, and a name can never come to mean different bytes.

What the file means accrues afterward, at leisure, as relationships. This scan is of that sample. This acquisition came out of that run. This image became figure 3. A person adds a link when they learn one, and so can a program. Nothing has to be right the first time, because a wrong link is superseded rather than repaired, and the wrong link stays, as a record that the mistake happened and when.

Save first, name by fingerprint, connect later.

What holds the connections is close to nothing. There are two tables. One holds objects, where an object is anything you want to refer to: a sample, a file, a person, an instrument, a run, a rule about how to read the other rules. The other holds relations, where a relation is a dated, labeled, confidence-weighted link between two objects, asserted by someone or something. This file was taken by that person, and I am fairly sure. There is no third table. There is no table for samples, and there is no table for measurements, and there is not going to be one.

Three commitments follow, and each will look wrong at first. Records are only ever added, so nothing is updated and nothing is deleted; a correction is a new record pointing at the one it supersedes, and reading the present means taking the latest word. The past is frozen by construction, which is what you want of a laboratory record and what almost no database gives you. A copy begins by writing down the moment it was switched on, and nothing in it precedes that moment, because an instance does not need a history it never had, only a beginning. That birth record is also how a copy knows which copy it is, since no two of them were born in the same instant. Identity is minted once and is never a description, so a thing's name is not what it is called and not what it is; names and categories and properties are links hung on an identity rather than the identity itself, and that is what lets you record something before deciding what it is. Entities are assembled rather than stored, so there is no master row for a given sample; a sample is the collection of records claimed to refer to the same physical object, and any description of it is produced on demand by gathering those records and summarizing them. Learn later that two labels were the same sample and you add a claim. You never merge rows, and you never lose the evidence of the merge.

None of this is a paper design. The same two tables carried a large and genuinely messy real dataset for a year: \labLineageCorpus{}. That was not the intended application. It was the stress test, and it is why the design can be described honestly instead of hopefully. Running something hard for a year tells you which of your ideas were real.

It has also left real bugs in place. This book names five of them and prints the command to reproduce each, and naming them that precisely turned out to be the thing that got four of them fixed. The view that computes the consensus of conflicting claims was reading the raw table instead of the filtered one, so nearly fifty thousand witnesses that had already been retracted were still voting; it reads the filtered view now, and that view is defined before it, so the ordering slip that caused it cannot recur. The function that stores a file by its fingerprint could be handed metadata that overwrote the fingerprint, which put a hole through the one invariant everything else rests on; the merge runs the other way around now, and the hash is checked. The identity column was a per-machine counter, so two copies each minted an object 1 and meant different things by it; it is a time-ordered UUID now, and two containers that have never communicated mint distinct names for distinct things and the same name for the same file.

So if you boot a copy and run those checks you will watch them pass, not fail. The argument for printing them survives that and is better for it, because the fastest way to decide whether to trust a document is still to run one of its claims; what has changed is that the claim is now a claim about a repair, and a repair is the harder thing to be honest about. So check what was not done. Nothing was deleted: every retracted row is still in the table, and all that happened is that the view stopped counting it as evidence. The correction is itself a row in the database, recording the defect, its cause, and the fix, which means you can read the repair instead of taking my word that it happened. A system whose whole thesis is that corrections are records rather than erasures does not get to fix its own bugs by quietly editing a file.

Two things the fixes do not buy. A view that reads the right table is not a law that is right: the consensus law is still a seed prior, the number in it was guessed and has never been fit to anything, and no amount of fixing the plumbing will fit it. And no two copies have ever exchanged a record, because the machinery for them to do so does not exist yet. The identity fix removes the reason merging could not have worked. That is not the same as showing that it does.

The fifth defect went a stranger way, and it is worth your suspicion rather than your applause. An ingestion had run twice and appended the same rule row twice. This book said flatly that both rows would be there forever, because nothing here is ever deleted, and that this was the design doing what it says rather than an embarrassment. Both rows are gone. They were not repaired. On \labLineageWipeDate{} I wiped the populated instance and started it over from an empty copy, and the duplicate went with everything else: \labLineageObjects{} objects, \labLineageRelations{} relations, the \labLineageScale{}, all of it. The old universe was not deleted either. It is a \labLineageSize{} file on a disk, frozen, readable, and no longer running. What replaced it holds ten records, which are nine rules and the note of its own birth.

Say plainly what that does and does not license. Starting over is the only operation here that removes history, and it works only because it is not a deletion inside the system. It is a different system with a new beginning, and the old one survives as an artifact that nobody can write to. Nothing is ever deleted stays true of every universe; it never said anything about how many universes there are. And the price is worth stating in the same breath as the result: the fix for one duplicated row was to lose \labLineageObjects{} objects, which was acceptable only because I judged that none of it was worth keeping, and because \labLineageSource{} it had all been derived from still exist and can be read in again. This is not a repair technique. It is available once, cheaply, at the beginning, and it gets expensive fast.

The system is named after a result in mathematics called the Yoneda lemma, and you do not need the mathematics to use it. Stripped of its machinery, the lemma says that a thing is completely determined by its relationships to everything else. There is nothing more to know about an object than the full set of ways it connects to the others. Read as an instruction for building things, that is the license for the inversion above. You never have to write down what a sample is. Record every relationship it takes part in, which wafer it was cut from, which run measured it, who fabricated it, which scan imaged it, which figure it ended up in, and you have captured it completely, and any description you want is those relationships gathered up when you ask. A conventional database asks you to define the thing first and hang the data off it. This one works in the direction data actually arrives.

Three kinds of reader, and the book is arranged so you can stop when you have what you came for. If you have data and want somewhere to put it, read Part I. You will boot your own empty copy, save a file so it cannot be quietly lost or altered, and connect it to the thing it came from. Part I assumes no mathematics and no knowledge of this project, every example runs against your own empty database, and every number in it is a query you can execute yourself. If you want to know why it is shaped this way, read Part II: two tables, three amendments, and a tower of rules that live in the database rather than in anyone's head. Read it when you are about to argue with the design. If you want to check the design rather than admire it, read Part III, which is where the thing came from. Part IV is a case study of the one instance that was run hard, including where it did not work.

One thing here is worth insisting on. Every non-obvious claim in this book is marked with its status: theorem, correspondence, analogy, or open. Each mark carries an obligation. A theorem needs a citation, and \tier{thm} says someone proved it. A correspondence needs a way to check it, and \tier{corr} names the check. An analogy has to say where it breaks, and \tier{anal} says where. An open question has to say what would settle it, and \tier{open} says what. A book that mixes proved results with suggestive analogies and does not say which is which is worse than useless, because a reader who catches one unmarked metaphor will rightly discount everything else, including the parts that are solid.

Which is the lemma pointed back at the book. A thing is known by its relationships, and so is this one. Every claim here is attached to something you can check, or it is marked as not yet attached to anything. If you find an unmarked claim, or a mark that flatters us, that is a bug, and the most useful thing you can do is demote one of our claims.
